\documentclass[../main.tex]{subfiles}
\begin{document}

\section{Placement over operator}
\label{sec:placement}

\begin{claim}{NEGATIVE}{F04}
The project began with a binary, dropout before against after the residual add, and replaced it with a family split, positions that feed a sublayer against positions on or beside the stream. If either carried the effect, it would survive a panel where the operator is held fixed and the strength is matched. Both are tested here and neither survives, and what remains is the evidence for the specific placement and for the axes that produced it.
\end{claim}

\subsection{The founding claim on the panel}

The founding claim was that dropout on each branch output, before the add, separates clean from triggered inputs better than dropout on the stream after it\finding{NEGATIVE}{F24}. \Cref{tab:pre-post} pairs the placements within cell over \PreMinusPostN{} cached cells. The gap is \PreMinusPostMatched{} at the matched rule with an interval of [\PreMinusPostMatchedLow{}, \PreMinusPostMatchedHigh{}] and \PreMinusPostAdaptive{} at the adaptive rule. The earlier reading of a large pre-residual advantage on CIFAR-10 came from sweeping the published placement outside its own rate window, which on \VitBlocks{} blocks sits an order of magnitude below the ConvNet grid, and it is the first of the conclusions in this project that a strength-matched comparison overturned. The refutation rests on the full panel, which is the evidence class of a CRITICAL claim, and it is the reason every comparison below is read at matched strength.

% generated by scripts/paper/tab_headline_checks.py from results/coverage/coverage.json, results/<folder>/psbd_metrics.json (71 cells) at 2026-09-23T11:29:35.303911+00:00, commit 07121b0ea68df7b530651d80d1c561600b445e34-dirty. Do not edit by hand.
\begin{table}[htbp]
\centering
\small
\caption{The founding question on the basis panel: dropout before each residual add against dropout after it, paired within model at the headline quantile, at the matched rule and at the adaptive rule.}
\label{tab:pre-post}
\begin{adjustbox}{max width=\linewidth}
\begin{tabular}{llrrl}
\toprule
comparison & rule & n & mean delta AUROC & 95\% CI \\
\midrule
pre\_residual minus post\_residual & adaptive & 69 & +0.012 & [-0.016, +0.040] \\
pre\_residual minus post\_residual & matched & 69 & +0.010 & [-0.024, +0.045] \\
\bottomrule
\end{tabular}
\end{adjustbox}
\end{table}


\subsection{The family split does not survive operator matching}

The claim that replaced it was that the effect lives in the family, sublayer inputs against residual-adjacent positions, measured on the earlier CIFAR-10 dropout panel with an interval excluding $0$\finding{NEGATIVE}{F04}. \Cref{tab:family-split} reads the same split on the basis panel with the operator held fixed at the token mask, the input-side placements against the residual-adjacent ones, each family's cell value the mean over its members. The gap is \FamilyGapMatched{} at the matched rule, interval [\FamilyGapMatchedLow{}, \FamilyGapMatchedHigh{}], and \FamilyGapAdaptive{} at the adaptive rule, interval [\FamilyGapAdaptiveLow{}, \FamilyGapAdaptiveHigh{}]. \Cref{tab:family-split-loo} refits it with each attack dropped and the smallest refit is \FamilyGapMinLeaveOneAttackOut{}. \Cref{tab:family-split-local-global} splits it by trigger locality, \FamilyGapLocal{} on local triggers and \FamilyGapGlobal{} on global ones, neither excluding $0$. The reading that mixes operators inside the families, \cref{tab:family-split-confounded} in the appendix, gives \FamilyGapConfoundedMatched{} with interval [\FamilyGapConfoundedMatchedLow{}, \FamilyGapConfoundedMatchedHigh{}], a sign reversal produced by the gaussian and scale-up members of the input-side family, which is what confounding position with operator does. The original effect was a dropout measurement on a single dataset, the basis panel carries no dropout at any input-side position, and the operator-matched dropout row of \cref{tab:family-split} prints \pending{} until that sweep runs. On the evidence in hand the family claim is withdrawn, and the runner-up placement of \cref{tab:basis-ranking} being a residual-adjacent token mask says the same thing from the other side.

% generated by scripts/paper/tab_family_split.py from results/coverage/coverage.json, configs/psbd_basis.json, results/<folder>/psbd_metrics.json (71 cells) at 2026-09-23T18:25:02.156153+00:00, commit 5a1f1c5814dc5676c093f3021490a849712a13f4-dirty. Do not edit by hand.
\begin{table}[htbp]
\centering
\small
\caption{The within-model family gap with the operator held fixed at token masking. Input-side (attention input, MLP input and both sublayer inputs) minus residual-adjacent (attention output before the add and stream after the attention add), each family's per-model value the mean AUROC at the headline quantile over whichever of its members are cached. Read at the matched 0.6 rate and separately at the adaptive 0.8 rate, 5000-resample bootstrap 95\% intervals.}
\label{tab:family-split}
\begin{adjustbox}{max width=\linewidth}
\begin{tabular}{lrrlrrl}
\toprule
subset & n match & gap matched06 & 95\% CI & n adapt & gap adaptive08 & 95\% CI \\
\midrule
all & 71 & +0.016 & [$-$0.001, +0.036] & 71 & +0.045 & [+0.016, +0.075] \\
CIFAR-10 & 18 & +0.027 & [$-$0.008, +0.062] & 18 & +0.005 & [$-$0.055, +0.068] \\
CIFAR-100 & 15 & $-$0.002 & [$-$0.030, +0.028] & 15 & +0.065 & [+0.018, +0.117] \\
GTSRB & 15 & +0.003 & [$-$0.029, +0.050] & 15 & +0.027 & [$-$0.009, +0.068] \\
Tiny ImageNet & 17 & $-$0.008 & [$-$0.026, +0.013] & 17 & +0.041 & [$-$0.013, +0.107] \\
SVHN & 3 & +0.126 & [$-$0.025, +0.239] & 3 & +0.122 & [$-$0.055, +0.267] \\
EuroSAT & 3 & +0.145 & [+0.100, +0.217] & 3 & +0.218 & [+0.033, +0.384] \\
rate 0.01 & 19 & +0.007 & [$-$0.021, +0.037] & 19 & +0.021 & [$-$0.015, +0.062] \\
rate 0.05 & 24 & +0.007 & [$-$0.018, +0.033] & 24 & +0.059 & [+0.025, +0.093] \\
rate 0.1 & 28 & +0.031 & [$-$0.002, +0.067] & 28 & +0.049 & [$-$0.011, +0.111] \\
hard attacks & 33 & +0.043 & [+0.013, +0.075] & 33 & +0.091 & [+0.038, +0.142] \\
easy attacks & 38 & $-$0.007 & [$-$0.025, +0.011] & 38 & +0.005 & [$-$0.017, +0.030] \\
\bottomrule
\end{tabular}
\end{adjustbox}
\end{table}

% generated by scripts/paper/tab_family_split.py from results/coverage/coverage.json, configs/psbd_basis.json, results/<folder>/psbd_metrics.json (65 cells) at 2026-09-23T16:59:49.157437+00:00, commit d4a8014354b8fcc8d014beec50fef60ed46fa47d-dirty. Do not edit by hand.
\begin{table}[htbp]
\centering
\small
\caption{The operator-matched gap at the matched 0.6 rate, refit with 1 attack dropped at a time over the 7 attacks present in the panel of 65 backdoored models, so a gap driven by a single attack would show as a refit crossing 0.}
\label{tab:family-split-loo}
\begin{adjustbox}{max width=\linewidth}
\begin{tabular}{lrrl}
\toprule
dropped attack & n & gap matched06 & 95\% CI \\
\midrule
without BadNets & 59 & +0.022 & [+0.003, +0.043] \\
without Blend & 57 & +0.017 & [$-$0.004, +0.039] \\
without BPP & 59 & +0.014 & [$-$0.006, +0.034] \\
without LF & 59 & +0.025 & [+0.005, +0.047] \\
without SIG & 66 & +0.006 & [$-$0.011, +0.023] \\
without TaCT & 60 & +0.020 & [+0.001, +0.041] \\
without WaNet & 66 & +0.013 & [$-$0.004, +0.031] \\
\bottomrule
\end{tabular}
\end{adjustbox}
\end{table}

% generated by scripts/paper/tab_family_split.py from results/coverage/coverage.json, configs/psbd_basis.json, results/<folder>/psbd_metrics.json (65 cells) at 2026-09-11T04:14:06.107762+00:00, commit 4a04e5dd204b14d8368fb7eeef7092bd331e00d9-dirty. Do not edit by hand.
\begin{table}[htbp]
\centering
\small
\caption{The family gap split by trigger locality, under the operator-matched token_mask definition and the confounded basis-tag definition. Local triggers occupy a fixed spatial patch (BadNets, Label-Consistent, TaCT), global triggers cover the whole image (Blend, BPP, LF, SIG, WaNet, Adaptive-Blend). Adaptive-Blend, Label-Consistent clear no model on this panel and contribute to neither row.}
\label{tab:family-split-local-global}
\begin{adjustbox}{max width=\linewidth}
\begin{tabular}{lrrlrrl}
\toprule
trigger locality & n match & gap matched06 & 95\% CI & n adapt & gap adaptive08 & 95\% CI \\
\midrule
local triggers, token\_mask & 23 & -0.009 & [-0.032, +0.015] & 23 & +0.048 & [+0.013, +0.082] \\
global triggers, token\_mask & 42 & +0.013 & [-0.006, +0.036] & 42 & +0.026 & [-0.013, +0.067] \\
local triggers, basis tags & 23 & +0.013 & [-0.020, +0.045] & 23 & +0.087 & [+0.049, +0.123] \\
global triggers, basis tags & 42 & -0.062 & [-0.078, -0.045] & 42 & -0.033 & [-0.057, -0.011] \\
\bottomrule
\end{tabular}
\end{adjustbox}
\end{table}


\subsection{What the axes do carry}

Holding the operator fixed and moving the site, \cref{tab:operators-deltas}, the token mask at the attention input beats the token mask at the MLP input by \AttentionInputMinusMlpInputTokenMask{} with interval [\AttentionInputMinusMlpInputTokenMaskLow{}, \AttentionInputMinusMlpInputTokenMaskHigh{}]\finding{STRONG}{F05}. Both sites are sublayer inputs and both are followed by a LayerNorm, so the site effect within the input-side family is larger than the family gap ever was, and it separates the sublayers rather than the families. Holding the site fixed and moving the operator, gaussian noise at the attention input loses to the token mask by \GaussianMinusTokenMaskAttentionNorm{}, interval [\GaussianMinusTokenMaskAttentionNormLow{}, \GaussianMinusTokenMaskAttentionNormHigh{}], and to the channel mask by \GaussianMinusChannelMaskAttentionNorm{}, while at the normalised MLP input the same noise beats the token mask by \GaussianMinusTokenMaskMlp{}, interval [\GaussianMinusTokenMaskMlpLow{}, \GaussianMinusTokenMaskMlpHigh{}]\finding{STRONG}{F06}. The sign of the operator effect depends on whether a LayerNorm follows the injection point, and \cref{sec:latent} gives the measured reason.

% generated by scripts/paper/tab_operators.py from results/coverage/coverage.json, results/<folder>/psbd_metrics.json (65 cells) at 2026-09-11T04:15:24.023920+00:00, commit 4a04e5dd204b14d8368fb7eeef7092bd331e00d9-dirty. Do not edit by hand.
\begin{table}[htbp]
\centering
\small
\caption{Mean AUROC at q0.25 for each placement of the operator and site comparison, at the matched 0.6 rate and the adaptive 0.8 rate, over the panel of 65 backdoored models. n is the number of models the placement covers under that rule.}
\label{tab:operators}
\begin{adjustbox}{max width=\linewidth}
\begin{tabular}{lllrrrr}
\toprule
placement & site & operator & n match & mean AUROC matched06 & n adapt & mean AUROC adaptive08 \\
\midrule
before\_attention\_norm\_token\_mask & before\_attention\_norm & token\_mask & 65 & 0.926 & 65 & 0.935 \\
before\_attention\_norm\_channel\_mask & before\_attention\_norm & channel\_mask & 65 & 0.832 & 65 & 0.880 \\
before\_attention\_norm\_gaussian & before\_attention\_norm & gaussian & 65 & 0.719 & 65 & 0.796 \\
before\_mlp\_gaussian & before\_mlp & gaussian & 65 & 0.873 & 65 & 0.868 \\
before\_mlp\_norm\_token\_mask & before\_mlp\_norm & token\_mask & 65 & 0.817 & 65 & 0.821 \\
\bottomrule
\end{tabular}
\end{adjustbox}
\end{table}

% generated by scripts/paper/tab_operators.py from results/coverage/coverage.json, results/<folder>/psbd_metrics.json (65 cells) at 2026-09-11T04:15:30.584459+00:00, commit 4a04e5dd204b14d8368fb7eeef7092bd331e00d9-dirty. Do not edit by hand.
\begin{table}[htbp]
\centering
\small
\caption{Paired AUROC deltas at the matched 0.6 rate, 5000-resample bootstrap 95\% intervals. A2 isolates the operator with the site held fixed, C1 isolates the site with the operator held fixed.}
\label{tab:operators-deltas}
\begin{adjustbox}{max width=\linewidth}
\begin{tabular}{lrrl}
\toprule
comparison & n & mean delta AUROC & 95\% CI \\
\midrule
A2: gaussian minus token\_mask at before\_attention\_norm & 65 & -0.206 & [-0.259, -0.156] \\
A2: gaussian minus channel\_mask at before\_attention\_norm & 65 & -0.113 & [-0.151, -0.078] \\
A2: gaussian minus token\_mask at before\_mlp & 65 & +0.056 & [+0.026, +0.087] \\
C1: before\_attention\_norm minus before\_mlp\_norm, both token\_mask & 65 & +0.109 & [+0.079, +0.138] \\
\bottomrule
\end{tabular}
\end{adjustbox}
\end{table}


The title's claim, that position dominates operator, has to be stated with its condition\finding{SUPPORTING}{F07}. On the basis panel the range of mean AUROC across the token mask's all-blocks positions is \PositionRangeTokenMaskMatched{} at the matched rule against an operator range of \OperatorRangeAttentionNormMatched{} at the attention input, a ratio of \PositionOverOperatorRatioMatched{}, and \PositionRangeTokenMaskAdaptive{} against \OperatorRangeAttentionNormAdaptive{} at the adaptive rule, a ratio of \PositionOverOperatorRatioAdaptive{}. The operator range is carried almost entirely by gaussian noise at the single position where a LayerNorm absorbs it. With the masking operators alone the operator range is \OperatorRangeAttentionNormMasksMatched{} and \OperatorRangeAttentionNormMasksAdaptive{}, and the ratio is \PositionOverOperatorRatioMasksMatched{} and \PositionOverOperatorRatioMasksAdaptive{}. Position dominates operator among operators that survive normalisation, and an additive operator placed before a norm is the exception that the mechanism chapter explains. An earlier ratio of this project had no runnable path back to a number, and these are the ones that do.

\subsection{Per-sample agreement across placements}

If every placement estimated the same per-sample quantity, their rankings of individual samples would agree. \Cref{tab:mech-operator-agreement-family} gives the mean Spearman correlation between placements' per-sample fractional PSU at their matched rates, \AgreementWithinInputSideClean{} within the input-side family and \AgreementAcrossFamiliesClean{} across families on clean samples, with the full matrix against the recommended placement in \cref{tab:mech-operator-agreement} in the appendix\finding{SUPPORTING}{F11}. The agreement between gaussian noise and the token mask at the same position is \AgreementGaussianVsTokenMaskClean{}, and between the IBD-PSC amplification and the recommended placement \AgreementGainScaleVsRecommendedClean{}. The placements agree about as much across families as within them, which is consistent with a single underlying quantity read through different Jacobians, and no more than moderately, which is why probe diversity has something to add in \cref{sec:limitations}.

% generated by scripts/paper/mech_operator_agreement.py from results/<folder>/psbd/<placement>/rate_*_{clean,backdoor}.pt at 2026-09-23T18:19:00.094983+00:00, commit 5a1f1c5814dc5676c093f3021490a849712a13f4-dirty. Do not edit by hand.
\begin{table}[htbp]
\centering
\small
\caption{B1: mean per-sample PSU agreement (Spearman rho, averaged over placement pairs and over models) within and across the input-side and residual-adjacent families, the axis H20 splits on.}
\label{tab:mech-operator-agreement-family}
\begin{adjustbox}{max width=\linewidth}
\begin{tabular}{lrrrr}
\toprule
comparison & clean rho & n pairs & backdoor rho & n pairs \\
\midrule
within input\_side & 0.472 & 55 & 0.264 & 55 \\
within residual\_adjacent & 0.582 & 28 & 0.349 & 28 \\
across families (input\_side vs residual\_adjacent) & 0.512 & 88 & 0.302 & 88 \\
\bottomrule
\end{tabular}
\end{adjustbox}
\end{table}


\end{document}
